\documentclass[11pt,a4paper]{article}
\usepackage[margin=14mm,top=12mm,bottom=12mm]{geometry}
\usepackage{fontspec}
\setmainfont{Satoshi}[Path=/Users/mlp/Library/Fonts/,Extension=.otf,UprightFont=*-Regular,BoldFont=*-Bold,ItalicFont=*-Italic,BoldItalicFont=*-BoldItalic]
\setmonofont[Scale=0.85]{JetBrainsMonoNL Nerd Font Mono}
\usepackage{tabularx,booktabs,array,enumitem,xcolor}
\usepackage{graphicx}
\usepackage[hidelinks]{hyperref}
\pagestyle{empty}
\setlength{\parindent}{0pt}
\setlength{\parskip}{3pt}
\setlist{nosep,leftmargin=*}
\definecolor{ink}{HTML}{18502A}  % moss, deck palette "Computation pad"
\newcommand{\hsec}[1]{\par\vspace{3pt}{\large\bfseries\color{ink}#1}\par\vspace{1pt}}
\begin{document}
{\LARGE\bfseries\color{ink} One loop, five versions}\\[1pt]
{\small IEEE SENSORS 2026 tutorial, Sunday 25 October: From Na\"ive Loops to FAIR Datasheets}

I used one simulated soft sensor, one staircase of presses and one recording loop. Each version fixes the bug the one before it left.

\hsec{What the five versions taught}
\renewcommand{\arraystretch}{1.12}
\begin{tabularx}{\linewidth}{@{}>{\bfseries}l >{\raggedright\arraybackslash}p{4.3cm} >{\raggedright\arraybackslash}X@{}}
\toprule
 & The loop & What it taught \\
\midrule
V0 & Move the stage, then sleep & Nothing goes into the file during the ramps or the holds. \\
V1 & Check the stage in a loop & The gaps are filled. Now disk stalls cause jitter (p99 about 60 ms). \\
V2 & Threads and a queue & Jitter is fixed. Each row gets its time when the loop reads it, so the time is late by 13.5 to 21.5 ms. If I press Start at a bad moment, it loses about a third of the frames, and its file cannot show it. \\
V3 & Time from the frame, plus a frame number & The lag drops to about 0. The file now shows the frames it lost. \\
V4 & The loop ticks at twice the frame rate & In the full-length run on the slides it lost 0\,\% at every start time. An earlier run lost up to 1.2\,\%, so timing still matters. If a frame is lost, the file shows it. \\
\bottomrule
\end{tabularx}

What the recording leaves out, no analysis can bring back. A passing test does not make the claim true.

\hsec{Run the demo yourself}
\textbf{On your laptop.} Get the code (link on the tutorial page, QR below). In the folder with \texttt{naive\_vs\_threaded.py}, run:

\begin{center}\texttt{uv run --with numpy --with matplotlib python naive\_vs\_threaded.py --fast --seed 2}\end{center}

\texttt{--fast} runs a shorter staircase than the full-length run on the slides. Add \texttt{--phase-sweep} to see the loss for start times from 0 to 18 ms.

\textbf{In Colab.} Open the notebook from the link on the tutorial page (QR below) and run the cells from top to bottom. There is nothing to install.

\hsec{Why your numbers will differ from the slides}
\begin{itemize}
\item The sensor and the disk are simulated. I assumed how the disk stalls, so the jitter numbers show what the problem looks like. They say nothing about your disk.
\item Your laptop, your operating system and its timer change when threads wake up. The same code gives you different jitter and different loss.
\item The start time matters most. V2 loses anywhere from 0 to 34.1\,\% as I move the start from 0 to 18 ms. Press Start a few milliseconds later and you get a different datasheet.
\item The slides show the full-length run. \texttt{--fast} is shorter, so the start and the end of the run weigh more in its numbers.
\end{itemize}
Compare V2 and V4 on your own machine, then write your two loss numbers below.

\hsec{Your run}
\renewcommand{\arraystretch}{1.9}
\begin{tabularx}{\linewidth}{@{}l X l X@{}}
Operating system: & \hrulefill & CPU: & \hrulefill \\
V2 frames lost (\%): & \hrulefill & V4 frames lost (\%): & \hrulefill \\
\end{tabularx}

\vspace{4pt}
\noindent\includegraphics[width=2.5cm]{../qr/landing.png}\hfill
\parbox[c]{\dimexpr\linewidth-3cm}{\textbf{Demo, handout and sources: scan the code or open this address}\\
\texttt{https://claude.ai/artifact/2mrHgNaJrZLdJKoNr26FrE}}
\end{document}
